\documentclass[a4paper,12pt,fleqn]{article}%fleqn pour aligner les equations à gauche
\usepackage{../../Styles/Cours}


\newcommand{\chnum}{Chapitre 1}
\newcommand{\chtitle}{Préparation d'un "liquide étonnant"}
\newcommand{\dtype}{TP 1}

\begin{document}
	\maketitle
	Dans un recueil de manipulations de chimie, on trouve le protocole expérimental permettant d'obtenir un "liquide étonnant". Pour le préparer, il faut prélever des quantités de matière. Comment prélever une quantité de matière d'une espèce chimique ?
	
	\begin{figure}[h!]
		\caption{Protocole expérimental pour préparer un "liquide étonnant"}
		\label{doc1}
		\textbf{Étape 1 : la réalisation de ce liquide étonnant nécessite la préparation de trois solutions A, B et C dont voici les compositions :}\\
		
		\begin{minipage}{.48\linewidth}
			\emph{Solution A :}\\
			
			\textit{Matériel} : fiole jaugée de \SI{50,0}{mL}, coupelle, sabot de pesée, balance de pesée à \SI{0,01}{g}, pissette d'eau distillée, bécher de \SI{100}{mL} noté \textbf{A}, pipette compte-goutte pour ajuster la fiole.\\
			\textit{Préparation} : Préparer \SI{50,0}{mL} de solution par dissolution de \SI{1,4E3}{\mole} d'acide ascorbique de formule brute \ce{C6H8O6}. Verser la solution dans un bécher noté \textbf{A}.
		\end{minipage}\hspace{.02\linewidth}
		\begin{minipage}{.48\linewidth}
			\emph{Solution B :}\\
			
			\textit{Matériel} : fiole jaugée de \SI{50,0}{mL}, bécher de \SI{100}{mL} noté \textbf{B}, éprouvette graduée de \SI{50,0}{mL}, pipette jaugée de \SI{5,0}{mL}.\\
			\textit{Produit} : solution \textbf{A} et solution de iodure de potassium à \SI{0,50}{\mole\per\liter}\\
			\textit{Préparation} : Par dilution,  \SI{50,0}{mL} d'une solution de iodure de potassium à \SI{0,10}{\mole\per\liter}.\\
			Dans un bécher de \SI{100}{mL} noté \textbf{B}, verser \SI{25,0}{mL} de solution d’iodure de potassium de concentration \SI{0,10}{\mole\per\liter} à l’aide d’une éprouvette graduée. Ajouter ensuite \SI{5,0}{mL} de solution \textbf{A}
			à l’aide d’une pipette jaugée
		\end{minipage}
		
		\begin{minipage}{.98\linewidth}
			\emph{Solution C :}\\
			\textit{Matériel} : bécher de \SI{100}{mL}, éprouvette de \SI{10,0}{mL}\\
			\textit{Produits} : peroxyde d'hydrogène (eau oxygénée de concentration $C = \SI{0,89}{\mole\per\liter}$), thiodène (empois
			d'amidon).\\
			\textit{Préparation} : Dans un bécher noté \textbf{C}, verser à l’aide d’une éprouvette graduée de \SI{10,0}{mL}, une quantité de matière
			de \SI{4,45E-3}{\mole} de peroxyde d’hydrogène, puis une quantité de matière de \SI{0,83}{\mole} d’eau. Introduire ensuite	une pointe de spatule de thiodène.
		\end{minipage}
		
		\vspace{2em}
		\textbf{Etape 2 : préparation du "liquide étonnant"}\\
		Mélanger dans un bécher de \SI{250}{mL} la totalité des solutions \textbf{B} et \textbf{C}.\\
		Patienter jusqu’à observer un changement.\\
		Ajouter \SI{5}{mL} de solution \textbf{A} et remuer. En l’absence de décoloration, renouveler l’ajout de \SI{5}{mL}. Patienter jusqu’à	observation d’un nouveau changement.
	\end{figure}
	
	\begin{figure}[h!]
		\caption{Données}
		\label{doc2}
		Masse volumique de l’eau à \SI{25}{\celsius} : $\rho(eau) = \SI{1,0}{\gram\per\milli\liter}$\\
		Masse molaire atomiques : $M(C) = \SI{12,0}{\gram\per\mole}$ ; $M(H)=\SI{1,0}{\gram\per\mole}$ ; $M(O) = \SI{16,0}{\gram\per\mole}$\\
		Formule semi-développée de l'acide ascorbique (solide pur) :
		\chemfig[atom sep = 2em]{HO-C(-[2]H)([6]-H)-C(-[2]H)([6]-OH)-CH*5(-C(-[6]OH)=C(-[7]OH)-C(=[1]O)-O-)}
	\end{figure}

\section*{Partie 1 : Calculs préliminaires}
Déterminer les masses et les volumes correspondants aux quantités de matières données dans le protocole. Vous donnerez à chaque fois la forme littérale, poserez les calculs et donnerez les résultats et les unités adaptées.\\
\underline{Pour la solution A :}\\

Acide ascorbique de formule brute \ce{C6H8O6} :\\

\vspace{3cm}
Eau :

\vspace{3cm}

\underline{Pour la solution B :}\\

Ion iodure \ce{I-} :

\vspace{3cm}

\underline{Pour la solution C :}\\
Peroxyde d'hydrogène \ce{H2O2} :
\vspace{3cm}

\section*{Partie 2 : Réalisation et interprétation}
Réaliser la préparation du "liquide étonnant"\\
En quoi ce liquide est-il étonnant ?

	
\end{document}